\documentclass[twocolumn]{article}

\usepackage{cmap} % для поиска в PDF
\usepackage[utf8]{inputenc}
\usepackage[T2A]{fontenc}
\usepackage[russian]{babel}
\usepackage{amsmath, amssymb, amsthm}
\usepackage{graphicx}
\usepackage{hyperref}

\title{GRA-Core: Иерархическая стабильность для субъект-ориентированных роев ИИ-агентов \\
\large{Оркестр поощряет локальную пену, обнуляя глобальную}}
\author{О. Разработчик}
\date{\today}

\begin{document}

\maketitle

\begin{abstract}
В данной работе представлен фреймворк GRA-Core, предназначенный для иерархической стабилизации роев ИИ-агентов с сохранением индивидуальности каждого участника. Основной принцип заключается в следующем: \emph{Оркестр поощряет пену ИИ-агентов, но обнуляет пену в оркестре}. Локальная пена (внутренние противоречия агента, его характер, «тараканы») разрешена и даже поощряется, в то время как глобальная пена (общая фальшь, конфликты, нарушение цели) строго обнуляется с помощью градиентной адаптации. Мы формализуем разделение этих понятий, приводим условия сходимости и демонстрируем работающую реализацию, включающую серверную часть API и интерактивную приборную панель.
\end{abstract}

\section{Введение}
Современные системы мультиагентного взаимодействия сталкиваются с фундаментальной дилеммой: стремление к глобальной согласованности часто подавляет разнообразие агентов. Предлагаемый подход решает эту проблему за счёт введения иерархического стабилизационного функционала, где каждый агент сохраняет ненулевое значение «локальной пены» $\Phi_i^{\text{local}}>0$, отражающей его характер, смещения и творческий шум, но оркестр в целом должен удовлетворять условию $\Phi_{\text{global}}=0$ (полное соответствие глобальной цели).

В основе работы лежит идея о том, что подавление индивидуальности не является необходимым условием для достижения коллективной эффективности. Напротив, локальная пена может быть источником креативности и адаптивности, тогда как глобальная пена — мерой дезорганизации и рассогласования.

\section{Формальная постановка}
Обозначим через $x$ глобальное состояние роя, а через $a_i$ и $s_i$ — состояние и личностные параметры агента $i$. Введём общий функционал:

\begin{equation}
J(x, s_{1..N}) = \Phi_{\text{global}}(x) + \sum_{i=1}^N \left( \alpha_i \Phi_i^{\text{local}}(a_i,s_i) + \beta_i \Phi_i^{\text{align}}(a_i,s_i,x) \right) + R,
\end{equation}

где:
\begin{itemize}
\item $\Phi_{\text{global}}(x)$ — глобальная пена (подлежит обнулению);
\item $\Phi_i^{\text{local}}$ — внутренняя пена агента (допускается, ограничена сверху константой $c_i$);
\item $\Phi_i^{\text{align}}$ — мера нарушения выравнивания агента с оркестром;
\item $R$ — регуляризационные члены.
\end{itemize}

Правило обновления GRA-Core имеет вид:

\begin{equation}
(x^{t+1}, s_{1..N}^{t+1}) = (x^t, s_{1..N}^t) - \eta \nabla J(x^t, s_{1..N}^t),
\end{equation}

с тем ограничением, что градиентные компоненты, соответствующие $\Phi_i^{\text{local}}$, демпфируются для сохранения условия $0 < \Phi_i^{\text{local}} \le c_i$.

\section{Условия стабильности}
Рой считается \emph{стабильным}, если выполнены следующие условия:

\begin{align}
\Phi_{\text{global}}(x^*) &= 0, \\
\Delta \Phi_{\text{global}}(x^*) &= 0, \\
\Delta^2 \Phi_{\text{global}}(x^*) &> 0, \\
0 &< \Phi_i^{\text{local}}(a_i^*, s_i^*) \le c_i \quad \forall i.
\end{align}

Первые три условия обеспечивают нахождение оркестра в локальном минимуме глобальной пены; последнее гарантирует, что каждый агент сохраняет свою индивидуальность. Условие $\Delta^2 \Phi_{\text{global}} > 0$ играет роль проверки устойчивости достигнутого минимума по аналогии с положительной определённостью второй производной.

\section{Архитектура реализации}
Предложенная реализация включает следующие компоненты:

\subsection{Серверная часть (FastAPI + GRA-Core)}
Бэкенд реализован на FastAPI и управляет агентами, вычисляет значения пены и выполняет шаги GRA. Основной класс \texttt{GRAOrchestra} предоставляет следующие методы:
\begin{itemize}
\item \texttt{add\_agent} — добавление агента с вектором характера;
\item \texttt{remove\_agent} — временное отключение агента (его пена сохраняется);
\item \texttt{step} — выполнение одного или нескольких шагов обнуления глобальной пены;
\item \texttt{\_recompute\_global\_pena} — пересчёт глобальной пены как суммы квадратов ошибок выравнивания и меры конфликта характеров.
\end{itemize}

Глобальная пена вычисляется по формуле:
\begin{equation}
\Phi_{\text{global}} = \sum_{i \in \text{connected}} (\text{align\_error}_i)^2 + 0.5 \cdot \text{conflict},
\end{equation}
где \texttt{conflict} — дисперсия первых двух компонент векторов характеров подключённых агентов.

API предоставляет следующие эндпоинты:
\begin{itemize}
\item \texttt{GET /state} — получение текущих значений пены, состояний агентов и флага стабильности;
\item \texttt{POST /step} — выполнение шагов обнуления GRA;
\item \texttt{POST /connect} — подключение ранее отсоединённого агента;
\item \texttt{POST /disconnect/\{id\}} — удаление агента из оркестра.
\end{itemize}

\subsection{Фронтенд-панель}
Интерактивная панель управления, написанная на HTML/CSS/JS, визуализирует текущие значения $\Phi_{\text{global}}$, $\Delta\Phi_{\text{global}}$, $\Delta^2\Phi_{\text{global}}$ и состояние стабильности. Панель также отображает список агентов с их индивидуальными значениями локальной пены и позволяет подключать/отключать агентов в реальном времени.

\section{Экспериментальная демонстрация}
Для демонстрации работы системы были использованы четыре предустановленных агента: \texttt{impressionist} (хаос 0.9, детализация 0.2), \texttt{paranoid} (подозрительность 0.8, точность 0.9), \texttt{disciplined} (скорость 0.6, качество 0.7) и \texttt{rebel} (антицентричность 0.9, креативность 0.8). Начальное состояние характеризуется ненулевой глобальной пеной.

В процессе выполнения GRA-шагов наблюдается монотонное убывание $\Phi_{\text{global}}$ при сохранении $\Phi_i^{\text{local}} > 0$ для всех агентов. Алгоритм шага включает:
\begin{itemize}
\item уменьшение ошибки выравнивания каждого агента: \texttt{align\_error} $= \max(0, \text{align\_error} - 0.1 \cdot \text{align\_error})$;
\item сохранение локальной пены с добавлением малого случайного возмущения: $\text{local\_pena} = \max(0.1, \min(\text{local\_pena} + \mathcal{N}(0, 0.02), 2.0))$.
\end{itemize}

Экспериментальные результаты подтверждают, что предложенный механизм позволяет достичь стабильного состояния, удовлетворяющего всем условиям стабильности.

\section{Заключение}
В работе представлен фреймворк GRA-Core для управления роями ИИ-агентов, основанный на принципе разделения локальной и глобальной пены. Предложена формальная математическая модель, включающая целевой функционал, правило обновления и условия стабильности. Разработана полная программная реализация, включающая серверную часть на FastAPI и веб-интерфейс для визуализации и управления. Экспериментальная демонстрация подтверждает работоспособность подхода.

Возможные направления будущих исследований включают:
\begin{itemize}
\item распространение фреймворка на большие языковые модели с использованием градиентной аппроксимации;
\item исследование влияния топологии связей между агентами на динамику стабилизации;
\item разработка адаптивных механизмов управления коэффициентами $\alpha_i$ и $\beta_i$ в зависимости от контекста.
\end{itemize}

\section*{Благодарности}
Автор выражает благодарность сообществу open-source за вдохновение и инструменты, использованные при реализации проекта. Исходный код и полная документация доступны в репозитории: \url{https://github.com/qqewq/gra-orchestra}.

\bibliographystyle{plain}
\begin{thebibliography}{1}

\bibitem{gra-orchestra}
GRA-Orchestra: Живой оркестр ИИ-агентов. \url{https://github.com/qqewq/gra-orchestra}.

\bibitem{fastapi}
F. Ram\'irez, ``FastAPI: Modern web framework for building APIs,'' 2018.

\bibitem{numpy}
C. R. Harris et al., ``Array programming with NumPy,'' Nature, vol. 585, pp. 357–362, 2020.

\end{thebibliography}

\end{document}